\section{EXECUTABLES REFERENCE}
\label{sec:executables}
This section describes each script individually. Section~\ref{sec:usage} already covers how they fit
together into the nightly workflow; here the focus is on each script's own inputs, outputs and notable
behaviour.

\subsection{bin/run.sh -- main nightly executable}
Entry point of a single night's simulation; invoked with no arguments, either by cron or manually
(Section~\ref{sec:usage}). It locates its own configuration via \texttt{conf.d/globals.var}, relative to
its own path, so it must be run from within \texttt{bin/}. On start it kills any leftover Meso-NH-related
processes from a previous, possibly stuck run (\texttt{pkill} on the executable names from
\texttt{globals.var}), rotates the log directory, then runs the full phase sequence described in
Section~\ref{sec:usage} (\texttt{prep\_list.sh}, \texttt{filedef.var}, PREP\_REAL, EXE,
\texttt{postproc.sh}, \texttt{plot\_python.sh}, compression, staging, \texttt{backup.sh}, notification,
hand-off). Exit status is 0 on success; any hard failure calls the \texttt{error} function (log + alert
mail + \texttt{exit 1}). Individual phases can be disabled with the \texttt{NOPREP}\slash\texttt{NORUN}\slash
\texttt{NODIAG}\slash\texttt{NOFICVAL}\slash\texttt{NOPLOT} debug switches in \texttt{globals.var}
(Section~\ref{sec:config}).

\subsection{bin/prep\_list.sh -- generate reference legacy-style scripts}
Invoked by \texttt{run.sh} as a separate sub-process (\texttt{env bash -c}) before the PREP\_REAL phase,
unless \texttt{PREPLISTTRUE=0}. Its header comment (``script used to generate runscripts with all the
namelists inside, Lascaux\&Masciadri style'') is accurate and matches what most of the script actually
does: it generates \textbf{two self-contained, standalone bash scripts}, written into \texttt{WORKDIR}:
\begin{lstlisting}
1.PREP_EXP_REAL_<version>_<site>_<N>MOD_<date>_<levels>lev_<tail>  : reproduces the
                                                    PREP_REAL/SPAWNING step
2.PREP_EXP_RUN_<version>_<site>_<N>MOD_<date>_<levels>lev_<tail>   : reproduces the
                                                    EXE (Meso-NH run) step
\end{lstlisting}
Each is fully self-contained -- exported shell variables, the relevant namelists embedded via heredocs,
the exact \texttt{mpirun} invocation -- written in the style of the pre-FAST Lascaux \& Masciadri
Meso-NH workflow scripts. They are meant as \textbf{reference material for a user unfamiliar with FAST}
who wants to see, or manually reproduce, what happens at the plain Meso-NH level; \texttt{run.sh}'s own
PREP\_REAL/EXE phases are separate, inline implementations and do \textbf{not} use these generated
scripts at all.

Generating them requires knowing which init files are available, so \texttt{prep\_list.sh} does source
\texttt{globals.var}, \texttt{times.var}, \texttt{run.var}, \texttt{filedef.var} and \texttt{backup.var}
(with \texttt{GENERATELIST=1}, which suppresses the verbose logging those files would otherwise
produce). Because it runs in a separate shell, none of the variables it resolves survive back into \texttt{run.sh}. At the end, \texttt{prep\_list.sh} deletes the REAL/RUN staging directories it used while generating the scripts (\texttt{rm -rf}); only the two generated scripts are left behind in \texttt{WORKDIR}.

If invoked interactively (with any command-line argument) it will instead prompt for\\
\texttt{YEAR}/\texttt{MONTH}/\texttt{DAY}/\texttt{DAYB} rather than reading them from \texttt{times.var}.

\subsection{bin/postproc.sh -- DIAG/FICVAL extraction}
Sourced (not executed as a sub-process) by \texttt{run.sh} right after the EXE phase, so it shares
\texttt{run.sh}'s environment. Unless \texttt{NODIAG=1}, it loops over every Meso-NH output
(\texttt{.lfi}) file, and for each nesting level listed in \texttt{DIAGLEVELARRAY} runs the \texttt{DIAG}
executable and \texttt{conv2dia}/\texttt{diaprog} for every combination of heights in
\texttt{DIAGINFHEIGHTS}/\texttt{DIAGSUPHEIGHTS} (\texttt{run.var}, Section~\ref{sec:config}), producing
the derived astroclimatic parameters. Unless \texttt{NOFICVAL=1}, it also runs the FICVAL extraction
(vertical profiles/cuts and 2D ground maps) for the keywords configured in
\texttt{EXTRACTLIST\_DOM\_FIRST}/\texttt{EXTRACTLIST\_DOM\_SECOND}/\texttt{PVLIST\_*}/\texttt{CVLIST\_*}.\\
Output goes to \texttt{DIAG\_MNH}/\texttt{ASCII\_DIR} (Section~\ref{sec:usage}).

\subsection{bin/plot\_python.sh -- plotting}
Sourced by \texttt{run.sh} after \texttt{postproc.sh}, unless \texttt{NOPLOT=1}. Drives the Python
scripts under \texttt{bin/conf.d/utils/Plot\_python/} (Section~\ref{sec:config}) to turn the ASCII
output of \texttt{postproc.sh} into the delivered figures under \texttt{PLOT\_DIR}, and computes
\texttt{WINDAVG} (the K=4 average wind speed) and the derived \texttt{WINDAVGHIGH} flag against
\texttt{WINDSPEEDHIGHLIM} (\texttt{run.var}), used by \texttt{backup.sh} as described in
Section~\ref{sec:usage}. Uses ImageMagick's \texttt{convert} for at least one figure format conversion
step.

It computes the full parameter suite on \texttt{DOM\_FIRST}: CN2 (both a ``new-algorithm''
wind-gradient-corrected product and, in parallel, the \texttt{\_OS18} alternative product, Section~
\ref{sec:config}), seeing, isoplanatic angle and coherence time, potential and absolute temperature,
relative humidity, precipitable water vapour, plus wind at several fixed levels and vertical profiles,
and 2D ground maps for most of the above. Among the wind products, it also computes statistics over
three fixed height windows tied to specific instrument requirements at the deployment site
(script-internal suffixes \texttt{\_linc\_low}\slash\texttt{\_linc\_high}, \texttt{\_flao},
\texttt{\_argos} -- currently 30--415\,m, 5072--9127\,m, 30--20000\,m and 30--560\,m respectively); edit
the ranges and labels directly in the script if a different deployment's instruments need different
windows. It unconditionally copies its ground-parameter results (seeing/wind/temperature/RH/PWV) into
\texttt{AUTOREGRESSION\_DIR} and \texttt{AUTOREGRESSION\_STORE\_DIR}, ready for the short-timescale
hand-off described in Section~\ref{sec:usage}, if that component is deployed.

\subsection{bin/backup.sh -- upload and local backup}
Sourced by \texttt{run.sh} at the very end of a successful run, unless \texttt{BACKUPLATER=1} (in which
case it must be invoked separately). Reads \texttt{backup.var} (Section~\ref{sec:config}) to decide, per
output category, whether to rsync it to the primary and backup remote servers
(\texttt{REMOTE\_HOST}\slash\texttt{REMOTE\_HOST\_BK}), and/or keep a local copy under
\texttt{BACKUP\_DIR} (or, for figures, \texttt{BACKUP\_FIGURES\_DIR}). When \texttt{UPLOADFIGS=0}
(\texttt{times.var})
it normally skips the figure upload entirely, \emph{except} that it still uploads the near-ground wind
figures/ASCII if \texttt{WINDAVGHIGH=1} (set by \texttt{plot\_python.sh}) -- the self-contained
high-wind exception described in Section~\ref{sec:usage}. Can also be run stand-alone
(\texttt{./backup.sh}) for a previously-staged run, in which case it detects it is not being sourced
from \texttt{run.sh} and sources \texttt{globals.var}, \texttt{times.var}, \texttt{run.var} and
\texttt{backup.var} itself first.

\subsection{bin/pgdgen.sh -- PGD (orography) generation}
Stand-alone script, not called by \texttt{run.sh}; run manually at installation time and again only if
the domain geometry or topography source files change (Section~\ref{sec:deps-install}). Sources
\texttt{globals.var} and \texttt{pgd.var}, checks/creates the working directories, then generates the
PGD files for every nesting level (outer domain plus each nested level, per \texttt{NESTING}) via the
internal \texttt{prep\_list\_pgd} function (which in turn drives Meso-NH's \texttt{PREP\_PGD}/
\texttt{PREP\_NEST\_PGD} executables using the \texttt{NLPGDFIRST}/\texttt{NLPGDNEST}/
\texttt{NLPGDPREPNEST} namelist templates). On completion it moves its own log files into
\texttt{PGDDIR}. Takes no arguments; progress and errors go to \texttt{LOGPGDFILE}.

\subsection{Benchmarking helpers (not part of nightly operation)}
\texttt{bin/Run\_lista.sh}, \texttt{bin/crossc2d.py} and \texttt{bin/multi.py} are used together to
replay a list of past initialization dates (read from a file named \texttt{LISTAgrib}) purely to measure
CPU/timing performance. \texttt{Run\_lista.sh} builds a \texttt{RUN\_ALL.sh} launcher from that list and
copies the two Python helpers into a scratch \texttt{CPUTEST/} directory. None of the three are invoked
anywhere in the operational workflow (cron or otherwise) and can be ignored for normal use.
